
\subsection{Installment 9: “Folkebladet” May 08, 1918}

Among us too there have been controversies, yes, schisms. Among us too words have been spoken that wear and wound. But over these words there reigns great disagreement, and there may come new schisms. These can be prevented—not by resorting to threats or the language of power, not by empty plotting or compromises—but by going the way of conviction. If there is agreement as to what the principles are and what they intend, then much will have been gained. It may be another question to decide whether the principles stand wholly and entirely on the ground of Scripture. First there must be agreement on what the thing is, before there can be discussion of its “how.”

That seeking to promote this understanding—and both Birkeland and I have sought it—has nothing to do with person and personal wrangling, is plain. It is small-minded thinking that writes, for example, Birkeland’s articles to the account of personal squabbling. Birkeland’s contributions have not concerned anything lying between persons, but something lying outside them, namely what the principles are and what they mean.

Peace comes with understanding. But it can also come in another way. It is possible to bury the principles and get a peace that is as silent as the grave’s. In the graveyard it is hushed and still: the silence of death. There one contends neither over personal issues nor over principles.

If the Free Church buries the principles—or whatever still remains of them—this will not, without more, mean the Free Church’s downfall. It can surely continue its existence, taking comfort in memories from former days, lowering its aim and writing “the enlivening and liberation of the individuals” instead of “the enlivening and liberation of the congregations.” But then it must also prepare itself to find an answer to the question: Why does S stand alone?

But even if the principles are allowed for some time yet to exist upon the surface of the earth, there is nevertheless danger that they may become a petrifaction, or a common platform from which everyone fetches what best suits his own need, whether it be a new Christian type, or an ideal, or something else.

Therefore: the way of conviction.

If the understanding of the principles were greater, the annual meetings would look different. As things now are, they are a kind of adding machine for the politics that are carried on between the annual meetings, an emergency brake to which the train staff resort every time they “see ghosts” on the line. So it was in Marinette, in Willmar, and in Fargo, where arrangements were even made to put some inconvenient passengers off. Had the understanding of the principles been greater, the meeting in Fargo would not have witnessed that concerted, peevish opposition to Birkeland and me. The meeting in Fargo was a classic expression of the situation among us.

Birkeland’s and my articles in Folkebladet took principle no. 2 as their point of departure. And this principle is, after all, acknowledged as essentially correct—even in the most conservative Lutheran circles in Europe. What our articles asserted were old Lutheran truths concerning the religious equality of the laity with the pastors. No attempt was made to overturn formal rights and arrangements, which were valued for precisely what they are—human arrangements.

I showed myself fully in agreement with Scripture and Lutheran doctrine and paid due regard to the principles. No one has refuted anything that B. and I wrote. The opponent’s weapons and shield were vaporous declamations about “spirit,” personalities, and “periphery.” Nothing tangible. Yet the loudest spokesmen against the ordination articles were certainly not the faithful followers in church law of the Coast pastor (J. J. Jansen), who himself also gives a detailed account of the problems of ordination. Jansen wrote—as we have earlier indicated—in his autobiography:

“Together with my aunts and my dreams of becoming a pastor, there had developed in me a sort of high-church enthusiasm for the church and the office, and with that also an intense interest in the literature as a whole. It was here that I meant to study, understand, and come to a result, first concerning what the New Testament teaches regarding the office. I read Ritschl, A. Rothe, Thiersch, Stahl, Löhe, Lechler, but on the other side also Höfling, Richter, Rudelbach. I studied Carpzov, Böhmer, Buseudorf, and many more. The great folios in my library still remind me of many an eager hour.....

“My heart was from the beginning with Stahl and Löhe, with the high-church, catholicizing, mystical doctrine of the church and doctrine of office. But in the course of study I became more and more low-church. How did I become low-church? The New Testament compelled me. Thus there came forth my, no doubt very imperfect, scientific attempt, ‘The Ecclesiastical Office in the Apostolic Age.’ (1877–78).”

It should be known that Jakob Sverdrup attacked this attempt in Ny luthersk Kirketidende. It was too high-church for Jakob Sverdrup, who himself wrote much about the office—low-churchly. But more low-church than Jakob Sverdrup was his brother Georg, who with respect to the office essentially shared Sohm’s view.

But when this low-church doctrine of office, with the corresponding contentious understanding of ordination, came forward again in Folkebladet, 1917, it seemed to some like something uncannily foreign and free-churchish. Our contributions were even perceived as personal attacks.

Birkeland’s and my articles were lumped together at the Fargo meeting. What the dissatisfied said about the one of us was also made to apply to the other. A remarkable, artificially constructed relation of solidarity! As for “personalities,” I knew myself quite free of them. And as regards Birkeland, it astonishes me that he could write as objectively as he did, after the unjust treatment given him in Willmar, and after the maneuvering that had been employed, and had succeeded, in preventing his reelection as chair of Augsburg’s Board of Trustees. It was said that money would not come in to Augsburg if he were to retain the chairman’s seat, and that confidence in the school would disappear! Even the ordination committees had brought Birkeland’s name into the questions they put to the theological candidates. In view of this and much else I say: all respect for Birkeland.

As for me, I have always sought to argue as one generally argues in respectable theological journals. It has been objective truth that has concerned me. When I have said or shown that this or that publicly advanced assertion is untenable, then I have said no more than a conscientious carpenter has the right and duty to say to another carpenter: This beam will not do. Get something else! But such an address is no personal attack. Nor is it murder—which, as I understand it, an editorial statement would underline. It is not for murder’s sake that holes are shot in untenable arguments. Coverings that serve to protect what is false—it is all one how well-intentioned they may be—deserve to be leveled.

My language has been the non-narcotizing language of duty. Such language speaks openly. It does not whisper, does not threaten. It does not whisper in the twilight, does not compose protests and does not circulate petitions. It does not maneuver under cover of committee chairmen’s reports. In its embarrassment it does not take refuge in words like “spirit and life,” words which in Scripture mean very much, but which here and there among us have in fact been degraded into a slogan that conceals kinds of diversity of many sorts, including intellectual irrationality and sloth.

The non-narcotizing language of duty has also been Birkeland’s language, yes, even where he cited Stendal’s words about Christ and Belial. For there are indeed collisions among us. And if one takes into consideration that the men who most often use the grand language about spirit and life took two days of the Fargo meeting’s costly time to discuss Birkeland’s (and my) contributions, then one gets the impression that the collisions are anything but child’s play. Yes, it seems as though there were a little (perhaps much) of a Christ-Belial relation of opposition here, where Birkeland and I become a kind of establishments of Belial in modern form, trafficking in personalities and needing publicly to be chastised and instructed concerning the necessity of the one thing needful, exactly as though we had never before heard or known what the first foundational truths of Christianity are.

I have always respected Birkeland for his labors on behalf of the principles. He is one of the few who has shown a firm grasp of the principles and has interpreted them consistently. If there is anything to object to in the consequences B. has set forth, then the fault is not Birkeland’s, but the principles’.

Now when Folkebladet has undertaken to humiliate Birkeland by editorial rebukes—boldly personal rebukes—and by publishing others’ disapproving criticism of his person without any sort of editorial dissent, it seems to me that he has good ground to take issue with Folkebladet, even if he sets aside the personal element in these articles and fastens upon the more substantive side.

1. Birkeland had good reason to criticize the decision that the Fargo meeting adopted in connection with the union matter. Neither the decision—which came as an overrumpling—nor the editorial defense, which was surely prepared in calm and ease, is in harmony with Leading Principles.

2. Birkeland also had good reason to criticize the editorial “Impressions from the annual meeting.” I must say I did not recognize the annual meeting in these impressions; nor did I know with what justification the editorial office expressed its hope that “those who hereafter write in the paper will respect the general sentiment among the people with regard to this matter.” But I did recognize the annual meeting in the criticism that P. Fostervold exercised upon the “Impressions.” Fostervold wrote among other things this:

“But how the paper can write such things is more than I can understand. When one tries to make it appear that the general sentiment among the people was on the side of those pastors who attacked the articles on ordination, then Folkebladet must know that it is doing the annual meeting an injustice. For every single layman who spoke in this matter—and they were not so few—placed himself clearly and decidedly against these pastors. And they surely, better than anyone else, represented the general sentiment among the common people. These men were not, either, some young brawlers. They were older, sober-minded men. There were, moreover, several pastors as well who wholly took the laity’s side in this matter. Besides, it was in the votes that the strongest words were spoken, and it was these that showed the general sentiment among the people in the meeting.”

So far as I was able to survey the situation at the annual meeting, two thirds of the meeting’s voting members cast their decisive votes on Birkeland’s side. The laity for a time thought of holding their own meeting and protesting against the dissatisfied pastors’ way of proceeding. But nothing came of it: they saw that it was unnecessary.

The chairman cast some light on the situation when, in characteristic fashion, he explained that the writings in Folkebladet (B.’s and mine) had widened the gulf between the pastors and the laity.

Unfortunately—there was indeed a gulf. It had come into being in former times. But a new element had been added: by these articles the people had, more than ever before, become conscious of this gulf. The gulf had been created long before—by the high-church doctrine of office, by the conception of the pastor’s special religious standing by virtue of call and ordination. For it was no secret that at a pastors’ meeting in Minneapolis it had been declared that a pastor stands in a nearer relation to God than a layman does.
